\section{Ход работы}

\subsection{Топология (пункт~\ref{task:topology})}
Собранная в симуляторе топология показана на рисунке~\ref{fig:topology}.

\bbtScreenshot{materials/topology.png}{Топология сети с именами интерфейсов}{fig:topology}

PC1 подключён к интерфейсу Gi0/0 маршрутизатора R1, PC2 — к Gi0/0 маршрутизатора R2.
Маршрутизаторы соединены интерфейсами Gi0/1. Под линиями связи подписаны сети, под
устройствами — последний байт их адреса.

\subsection{Адресация (пункт~\ref{task:address})}
Адреса назначены по плану из таблицы~\ref{tab:plan} и сведены в таблицу~\ref{tab:addresses}.

\begin{bbttable}{Адреса интерфейсов}{tab:addresses}
  \begin{tabular}{|l|l|l|l|l|}
    \hline
    Устройство & Интерфейс & IP-адрес & Маска & Шлюз \\ \hline
    PC1 & FastEthernet0 & 192.168.5.10 & 255.255.255.0 & 192.168.5.1 \\ \hline
    R1 & Gi0/0 & 192.168.5.1 & 255.255.255.0 & — \\ \hline
    R1 & Gi0/1 & 10.0.5.1 & 255.255.255.252 & — \\ \hline
    R2 & Gi0/1 & 10.0.5.2 & 255.255.255.252 & — \\ \hline
    R2 & Gi0/0 & 192.168.15.1 & 255.255.255.0 & — \\ \hline
    PC2 & FastEthernet0 & 192.168.15.10 & 255.255.255.0 & 192.168.15.1 \\ \hline
  \end{tabular}
\end{bbttable}

Команды настройки интерфейсов R1 приведены в листинге~\ref{lst:r1}. На R2 введены те же
команды со своими адресами из таблицы~\ref{tab:addresses}.

\bbtCodeFile[text, lines=2-14]{materials/r1-config.txt}{Настройка интерфейсов
R1}{lst:r1}

Каждый интерфейс получает адрес и маску командой \texttt{ip address} и включается командой
\texttt{no shutdown}: по умолчанию интерфейсы маршрутизатора выключены. Маска
255.255.255.252 оставляет в сети между маршрутизаторами ровно два адреса узлов — для R1 и
R2.

\subsection{Связь до настройки маршрутов (пункт~\ref{task:before})}
Результат команды \texttt{ping} с PC1 на PC2 до настройки маршрутов показан на
рисунке~\ref{fig:before}.

\bbtScreenshot[0.85\linewidth]{materials/ping-before.png}{Проверка связи PC1 с PC2 до
настройки маршрутов}{fig:before}

На все запросы отвечает не PC2, а R1 с адреса 192.168.5.1: «Destination host unreachable».
R1 знает только свои подключённые сети, маршрута в сеть 192.168.15.0/24 у него нет, поэтому
он отбрасывает пакеты.

\subsection{Статические маршруты (пункт~\ref{task:routes})}
На R1 добавлен маршрут в LAN2 через R2 (листинг~\ref{lst:route1}), на R2 — маршрут в LAN1
через R1 (листинг~\ref{lst:route2}).

\bbtCodeFile[text, lines=15]{materials/r1-config.txt}{Статический маршрут на R1}{lst:route1}

\bbtCodeFile[text, lines=15]{materials/r2-config.txt}{Статический маршрут на R2}{lst:route2}

Маршрут нужен на обоих маршрутизаторах: R1 доставит запрос PC1 до LAN2, но ответ PC2 пойдёт
обратно через R2, и без своего маршрута R2 отбросит его. Таблица маршрутизации R1 после
настройки показана на рисунке~\ref{fig:routes}.

\bbtScreenshot[0.85\linewidth]{materials/r1-routes.png}{Таблица маршрутизации
R1}{fig:routes}

В таблице две подключённые сети с кодом C — 10.0.5.0/30 и 192.168.5.0/24 — и статический
маршрут с кодом S в сеть 192.168.15.0/24 через 10.0.5.2, то есть через R2.

\subsection{Проверка связи (пункт~\ref{task:check})}
Результаты команд \texttt{ping} и \texttt{tracert} с PC1 на PC2 после настройки маршрутов
показаны на рисунке~\ref{fig:after}.

\bbtScreenshot[0.85\linewidth]{materials/ping-after.png}{Проверка связи PC1 с PC2 после
настройки маршрутов}{fig:after}

На все четыре запроса ответил сам PC2, потерь нет. Команда \texttt{tracert} показывает путь
пакета: шлюз PC1 — R1 (192.168.5.1), затем R2 (10.0.5.2) и сам PC2 (192.168.15.10). Путь
совпадает со статическими маршрутами из листингов~\ref{lst:route1} и~\ref{lst:route2}.
